% Experimental setup

\section{Experimental setup}
\label{sec:setup}

The validation follows the reviewer's acceptance ladder
(CHECKLIST, items L1--L4): a laminar smoke test, forced turbulent runs at
three Reynolds numbers, a full-grid spectral reference, and a static POD
baseline.

\subsection{L1: Taylor--Green laminar decay}
\label{sec:taylor-green}

The unforced problem ($F = 0$) on $\T^2$, initialized with the Taylor--Green
stream function
\begin{equation}
  \psi_{\mathrm{TG}}(x,y,0) = A \sin x \sin y
  \qquad\Longleftrightarrow\qquad
  u_{\mathrm{TG}} = (A \sin x \cos y,\; -A \cos x \sin y),
  \label{eq:tg}
\end{equation}
with $A = 1$
% [PENDING-CODER: confirm amplitude A, viscosity nu (Re), grid N, time span,
% initial rank r_0, and rank tolerances for the L1 run.]
and viscosity $\nu$ fixed by the chosen Reynolds number
$\mathrm{Re} = 2\pi/\nu^2$
% [PENDING-CODER: confirm Re value used for the laminar decay run.]
, the flow decays monotonically: $E(t)$ is non-increasing by
\eqref{eq:energy}, and the singular values of $\psi(\cdot,\cdot,t)$ decay
rapidly. This test verifies (i) the discrete energy monotonicity (I2,
unforced), (ii) the exact divergence-free property (I1), and (iii) the
$\sigma_{\mathrm{tol}}$ decay rule of Section~\ref{sec:rank-adaptation},
which should drive the rank $3 \to 2 \to 1$ as the vortex dissipates
(Figure~\ref{fig:tg}).

\subsection{L2: forced turbulent runs}
\label{sec:forced}

The main experiments are forced Kolmogorov-flow runs with
$f = F\sin y\, e_x$ at
\begin{equation}
  \mathrm{Re} \in \{100,\; 1000,\; 5000\},
  \qquad \mathrm{Re} = 2\pi F/\nu^2,
  \label{eq:re-set}
\end{equation}
% [PENDING-CODER: confirm the (F, nu) pairing for each Re in
% {100, 1000, 5000}.]
on grid resolution $N = \text{[PENDING-CODER]}$
% [PENDING-CODER: grid resolution N (and any resolution dependence check) for
% each Re.]
with time step $\Delta t = \text{[PENDING-CODER]}$
% [PENDING-CODER: time step / time-stepping policy for each Re.]
and initial condition
$\text{[PENDING-CODER: initial condition, e.g.\ Kolmogorov equilibrium plus
perturbation]}$.
% [PENDING-CODER: confirm IC and total simulation time, including the length
% of the spin-up (discarded) phase and the statistical (retained) window.]
Each run records the invariants I1--I4 of
Section~\ref{sec:invariants}: rank over time, singular-value decay,
divergence-free residuals, kinetic energy and its spectrum, and per-step
costs.

\subsection{L3: full-grid spectral reference}
\label{sec:reference}

For accuracy and invariant comparison we integrate the same
stream-function formulation \eqref{eq:stream-ns} on the full grid with a
spectral method (Fourier collocation, exact $\Delta$ and $\Delta^{-1}$,
same de-aliasing policy and time-stepping policy as the SP-DLRA runs)
\cite{orszag1971}.
% [PENDING-CODER: confirm the reference scheme's integrator, de-aliasing,
% time step, and grid (possibly finer than the SP-DLRA grid), and how
% reference fields are sampled for error/invariant comparison.]
The reference solution $\psi_{\mathrm{ref}}$ is the ground truth for the
relative error
$\norm{\Psi - \psi_{\mathrm{ref}}}_2 / \norm{\psi_{\mathrm{ref}}}_2$ and for
the kinetic-energy statistics of I4.

\subsection{L4: static POD baseline}
\label{sec:pod}

As the static-basis baseline we build an offline proper orthogonal
decomposition (POD) from snapshots of the reference run of each
Reynolds number \cite{sirovich1987, lumley1967}: the snapshot matrix is
eigen-decomposed and truncated at the smallest number of modes
$r_{\mathrm{POD}}$ resolving 99.9\% of the (snapshot) energy. The baseline
model is the Galerkin projection of the dynamics onto this \emph{fixed}
basis (no online rank adaptation), advanced with the same nonlinear
residual evaluation as SP-DLRA.
% [PENDING-CODER: confirm the snapshot window used for the POD (statistical
% window only), the exact energy threshold (99.9%), the Galerkin baseline
% integrator, and how the POD basis handles the forcing term.]
The POD baseline measures the cost of staticity: whether a fixed basis of
size $r_{\mathrm{POD}}$ can track the dynamics at a cost and accuracy
comparable to the adaptive rank-$r(t)$ SP-DLRA model.

\subsection{Metrics}
\label{sec:metrics}

We report, per Reynolds number:
\begin{itemize}
  \item \emph{Rank dynamics}: $r(t)$ and the singular-value spectrum
  $\sigma_i(t)$ (Figures~\ref{fig:rank} and~\ref{fig:svd});
  \item \emph{Accuracy}: $\max$ relative $L^2$ error against the
  full-grid reference over the statistical window, per Re
  (Figure~\ref{fig:error});
  \item \emph{Cost}: per-step and total wall-clock time and peak memory for
  SP-DLRA, the full-grid reference, and the POD baseline, including the
  regimes in which SP-DLRA is slower (Figure~\ref{fig:cost},
  Section~\ref{sec:cost});
  \item \emph{Invariants}: $\max|\grad\cdot u|$ per run (Table~\ref{tab:div}),
  kinetic-energy monotonicity (I2), forcing-aware energy balance (I2), rank
  versus $r_{\mathrm{POD}}$ (I3), and kinetic-energy spectrum / statistics
  against the reference (I4).
\end{itemize}

\subsection{Hardware and reproducibility}
\label{sec:repro}

All experiments are single-node.
% [PENDING-CODER: hardware (CPU/GPU), software stack and versions, and total
% compute budget per (Re, N).]
Per the project's reproducibility checklist (CHECKLIST, item 1.1), every
number and figure in the paper is traceable to a run record in
\texttt{state/coder/results/} with its configuration, driver script, and
source commit; figures are generated from those records by scripts under
\texttt{experiments/}.
% [PENDING-CODER: confirm the results-directory layout and the figure
% generation scripts so the provenance note above is accurate.]
